\subsection{Representational mechanism and testable predictions}
\label{sec:hdc-theory}

Hyperdimensional computing (HDC) represents values with fixed distributed
codes and combines them through binding and bundling
\citep{kanerva2009,kleyko2023survey}. Here, the front end encodes each neural
channel's bounded continuous value, binds the value code to a channel role,
and superposes the channel contributions. It then binds recent population
codes to temporal roles before a causal recurrent decoder. The resulting
similarity kernel and perturbation response motivate the representation
comparisons below.

\subsubsection{Channel-value binding}

Let $x_{t,c}$ be the normalized, causally smoothed value of channel $c$ at bin
$t$. The encoder maps it to $u_{t,c}=(1+\tanh x_{t,c})/2\in(0,1)$ and linearly
interpolates between the two neighboring points on a fixed grid of $L$ anchors.
Write the resulting nonnegative weights as $a_{t,c,\ell}$, with
\[
  \sum_{\ell=1}^{L} a_{t,c,\ell}=1,
  \qquad
  a_{t,c,\ell}=0\ \text{except at the two adjacent anchors.}
\]
Each anchor $\ell$ has a fixed complex unit code
$v_\ell=\exp(i\,a_\ell\omega)\in\mathbb{C}^D$, where $a_\ell$ is its
normalized grid position and each coordinate of $\omega$ is drawn
independently from a zero-mean Gaussian with standard deviation $2\pi$. Each
channel $c$ has a fixed complex unit role $r_c\in\mathbb{C}^D$. Their
elementwise complex product binds value and channel identity. The population
hypervector is
\[
  q_t=\frac{1}{\sqrt{C}}\sum_{c=1}^{C}\sum_{\ell=1}^{L}
       a_{t,c,\ell}\,(r_c\odot v_\ell)\in\mathbb{C}^{D}.
\]
Complex representations are supplied to learned layers as concatenated real
and imaginary coordinates, retaining input-dependent phase. The codebooks are reproducible functions of
$\texttt{codebook\_seed}$, which is deterministically paired with the training
seed. Each final replicate therefore varies both codebook and optimization
randomness; the two sources of variability are not estimated separately.

The implementation flattens channel-anchor weights into
$\alpha_t\in\mathbb{R}^{CL}$ and precomputes the bound basis matrix
$M\in\mathbb{R}^{CL\times 2D}$. It evaluates $\alpha_tM$ with one matrix
multiplication. The materialized rate encoding has shape
$N\times T\times C\times L$, where $N$ is batch size, rather than an
$N\times T\times C\times D$ activation.

For independently sampled uniform channel roles and Gaussian value
frequencies, define normalized similarity as the real part of the complex
inner product divided by hyperdimension:
\[
 \begin{aligned}
 k(x,x')&=\mathbb{E}_{r,\omega}\!\left[\frac{1}{D}
                    \operatorname{Re}\{q(x)^Hq(x')\}\right]\\
   &=\frac{1}{C}\sum_{c=1}^{C}\sum_{\ell=1}^{L}
      \sum_{m=1}^{L}a_{c,\ell}(x)\,a_{c,m}(x')
      \exp\!\left[-\frac{(2\pi)^2}{2}(a_\ell-a_m)^2\right].
 \end{aligned}
\]
Distinct channel roles have zero expected cross-product, while the Gaussian
characteristic function gives the anchor-value similarity kernel. The anchor
vectors use a random Fourier construction \citep{rahimi2007features}; linear
interpolation extends this finite codebook to continuous values and yields a
finite-grid approximation to graded Gaussian similarity. Channel binding
combines per-channel similarities in a shared representation, with channel
identity remaining explicit through the distinct role vectors.

For a single bounded value $u\in[0,1]$, its interpolation weight at the
uniform anchor position $a_\ell=(\ell-1)/(L-1)$ is
\[
  p_\ell(u)=\max\!\left\{0,1-\left|(L-1)u-(\ell-1)\right|\right\}.
\]
The exact finite-grid kernel for one channel is
\[
  \mathcal{K}_L(u,u')=\sum_{\ell=1}^{L}\sum_{m=1}^{L}
      p_\ell(u)p_m(u')
      \exp\!\left[-\frac{(2\pi)^2}{2}(a_\ell-a_m)^2\right].
\]
Because the anchor grid is finite, this kernel can depend on the absolute
values as well as their separation. The curve in Fig.~\ref{fig:architecture}
therefore averages the kernel uniformly over valid value pairs with a fixed
separation $\delta$:
\[
  \kappa_L(\delta)=\frac{1}{1-\delta}
      \int_{0}^{1-\delta}\mathcal{K}_L(u,u+\delta)\,du,
      \qquad 0\leq\delta<1,
\]
with the endpoint at $\delta=1$ given by the continuous limit. The plotted
quantity is $\kappa_8(\delta)/\kappa_8(0)$.

The expected squared normalized population perturbation has a related exact
form. For two inputs $x,x'$ and expectation only over independently sampled
channel roles, with the value codebook held fixed,
\[
 \mathbb{E}_{r}\!\left[\frac{\|q(x)-q(x')\|_2^2}{D}\right]
   =\frac{1}{C}\sum_{c=1}^{C}
       \frac{\|v(x_c)-v(x'_c)\|_2^2}{D}.
\]
Here $v(x_c)=\sum_\ell a_{c,\ell}(x)v_\ell$. Since interpolation is convex
and each anchor vector has norm $\sqrt{D}$,
$\|v(x_c)\|_2\leq\sqrt{D}$. If the inputs differ in at most $m$ channels,
the expression is bounded by $4m/C$. This bounds expected representation
change under channel-role randomness; it is not a prediction-accuracy
guarantee or a guarantee of invariance to channel-identity changes.

The fixed random-feature control receives the same interpolated
$CL$-dimensional input and uses an independent Gaussian mixing matrix with the
same output width. Both fixed encoders have no trainable parameters. At a
matched $(D,K)$ candidate, HDC and random have identical downstream learned
architectures and total trainable parameter counts. Their first learned stage
is a linear projection from $2D$ to the common hidden width, followed by layer
normalization. The random matrix's expected feature inner
product is proportional to the dot product of interpolated channel-anchor
weights. If $\alpha(x)\in\mathbb{R}^{CL}$ is the flattened interpolated input
and each entry of the random matrix $G$ has variance $1/(2C)$, then
\[
  \mathbb{E}_{G}\!\left[\frac{1}{D}
    \langle \alpha(x)G,\alpha(x')G\rangle\right]
  =\frac{1}{C}\langle\alpha(x),\alpha(x')\rangle.
\]
The HDC codebook instead induces the Gaussian anchor-similarity kernel above.
For the random control, the single-channel kernel is
$\mathcal{K}^{\mathrm{random}}_L(u,u')=\sum_\ell p_\ell(u)p_\ell(u')$.
Figure~\ref{fig:architecture} applies the same pair averaging to both
kernels and normalizes each by its own value at zero separation. The random
kernel vanishes when the two interpolation supports do not overlap; HDC's
correlated anchors retain graded similarity across a wider value range.
This difference specifies the feature geometry tested by the matched
random-control comparison.

The MLP control receives the normalized neural input as a $K$-bin lag stack and
applies a trainable nonlinear projection to the common hidden width. Its
bottleneck width is
\[
 w=\operatorname{round}\!\left(\frac{2D\,H}{C K+H+1}\right),
\]
where $H$ is the recurrent hidden size. This makes the MLP projection's
trainable parameter count within 1\% of the HDC/random projection for the
searched dimensions and context lengths. The plain control uses a learned
linear projection of the same $C K$ normalized lag stack. Every family searches the
same context lengths, $K\in\{1,8\}$, but the selected context length can differ
between families. The fixed mappings add no information about the target
beyond the input; either mapping may help a finite-capacity downstream decoder
use that input, which is an empirical question.

\subsubsection{Causal temporal role binding}

For a temporal window of $K$ bins, each lag $j$ has a fixed complex unit role
$\rho_j\in\mathbb{C}^D$. The temporal representation is
\[
  z_t=\frac{1}{\sqrt{K}}\sum_{j=0}^{K-1}
        \rho_j\odot q_{t-j}.
\]
Complex multiplication binds a population code to its lag role; addition
bundles the ordered recent codes. For independent random temporal roles, the
expected similarity between two bundled windows retains aligned-lag
similarities, while cross-lag terms have zero expectation:
\[
 k_{\mathrm{temp}}(t,s)
  =\mathbb{E}_{\rho}\!\left[\frac{1}{D}
          \operatorname{Re}\{z_t^Hz'_s\}\right]
   =\frac{1}{K}\sum_{j=0}^{K-1}
       \frac{1}{D}\operatorname{Re}\{q_{t-j}^Hq'_{s-j}\}.
\]
This gives an order-sensitive hypothesis: matching activity at the same
relative lags contributes to similarity. For a given candidate value of $K$,
the HDC, random, MLP, and plain families receive the same causal context
before the common GRU. The selected context length can differ across families.
The \texttt{no\_roles} ablation replaces each temporal role with the all-one
vector while preserving the lag window and bundling operation. At $K=1$, this
changes only a fixed coordinate rotation and does not test temporal ordering. The
\texttt{no\_channel\_roles} ablation replaces all $r_c$ with the all-one
vector, removing explicit channel identity from HDC value binding.

The streaming state stores the preceding $K-1$ encoded population vectors for
HDC/random, or the preceding $K-1$ normalized neural vectors for MLP/plain. Each
candidate uses only current and preceding bins. Chunked inference is
equivalent to single-pass inference up to floating-point summation order. State
is local to the decoder instance and is cleared at recording reset.

\subsubsection{Cross-session rationale and testable comparisons}

The hypothesis is that channel-value codes and explicit lag roles may provide
a bounded feature basis for a readout trained on other sessions. The bounded
value transform limits the effect of extreme normalized rates, while temporal
binding presents short ordered patterns to the recurrent decoder. Their
contribution is tested by comparing the selected HDC configuration with the
same-basis random projection, learned MLP, and conditional channel-role and
temporal-role ablations on held-out sessions. Synthetic gain and channel
dropout tests probe sensitivity to those specified input perturbations.

The representation does not align neural manifolds or identify corresponding
units across sessions. A fixed channel role assumes that channel index $c$
retains its meaning within the input schema; changes in unit identity or
channel ordering alter the encoded vector. Any observed transfer therefore
applies to the specified normalization and channel mapping. The evaluation
holds these choices fixed across representation families.
